%%%%%%%%%%%%%%%%%%%%%%%%%%%%%%%%%%%%%%%%%%%%%%%%%%%%%%%%%%%%
% (01) Input Settings
%%%%%%%%%%%%%%%%%%%%%%%%%%%%%%%%%%%%%%%%%%%%%%%%%%%%%%%%%%%%
\input{settings.tex}

%%%%%%%%%%%%%%%%%%%%%%%%%%%%%%%%%%%%%%%%%%%%%%%%%%%%%%%%%%%%
% (02) DOCUMENT INFORMATION - FROM YAML METADATA
%%%%%%%%%%%%%%%%%%%%%%%%%%%%%%%%%%%%%%%%%%%%%%%%%%%%%%%%%%%%
\newcommand{\DocTitle}{شی‌گرایی (ارث‌بری)}

\newcommand{\CourseTitle}{برنامه‌سازی پیشرفته و کارگاه}

\newcommand{\Semester}{بهار ۱۴۰۵}

\newcommand{\FirstProfessor}{دکتر روح‌الله احمدیان}

\newcommand{\SecondProfessor}{}

\newcommand{\Writers}{سام قربانی، آرمان حسینی، محمدحسین هاشمی، کیانا
پهلوان، کیانا رضوی، اریا فیض‌ بخش، مهدی جعفری، مارال مهاجر}

\newcommand{\Faculty}{دانشکده ریاضی و علوم کامپیوتر}

\newcommand{\University}{دانشگاه صنعتی امیرکبیر}

%%%%%%%%%%%%%%%%%%%%%%%%%%%%%%%%%%%%%%%%%%%%%%%%%%%%%%%%%%%%
% DOCUMENT BODY
%%%%%%%%%%%%%%%%%%%%%%%%%%%%%%%%%%%%%%%%%%%%%%%%%%%%%%%%%%%%
\begin{document}
	
%%%%%%%%%%%%%%%%%%%%%%%%%%%%%%%%%%%%%%%%%%%%%%%%%%%%%%%%%%%%
% === TITLE PAGE ===
%%%%%%%%%%%%%%%%%%%%%%%%%%%%%%%%%%%%%%%%%%%%%%%%%%%%%%%%%%%%
\begin{titlepage}
	\centering
	\vspace*{1cm}
	
	%— Smaller logo (20% of text width)
	\begin{figure}[H]
		\centering
		
\includegraphics[width=0.20\textwidth]{images/logo-fa.png}
	\end{figure}
	
	\vspace{1cm}
	
	{\fontsize{18}{28}\selectfont \textbf{\CourseTitle} }
	
	\vspace{3cm}
	
	{\fontsize{22}{28}\selectfont \textbf{\DocTitle} }
	
	\vspace{3cm}
	
	{\fontsize{14}{22}\selectfont \rl{نگارش:}}\\
	\vspace{0.5cm}
	{\fontsize{16}{22}\selectfont \textbf{ \Writers } }

	
	\vspace{1cm}
	
	{\fontsize{14}{22}\selectfont \rl{استاد درس:}}\\
	\vspace{0.5cm}
	{\fontsize{16}{22}\selectfont \textbf{ \FirstProfessor } }
	
	
	\vfill
	
	{\fontsize{14}{18}\selectfont \textbf{\Faculty} }\\
	\vspace{0.5cm}
	{\fontsize{14}{18}\selectfont \textbf{\University} }
	\vspace{0.5cm}
	
	\par
	
	{\fontsize{14}{18}\selectfont \textbf{\Semester} }
	
\end{titlepage}

%— Switch to our fancy header/footer from this page onward
\clearpage

%%%%%%%%%%%%%%%%%%%%%%%%%%%%%%%%%%%%%%%%%%%%%%%%%%%%%%%%%%%%
% === TABLE OF CONTENTS ===
%%%%%%%%%%%%%%%%%%%%%%%%%%%%%%%%%%%%%%%%%%%%%%%%%%%%%%%%%%%%
\tableofcontents
\clearpage
\pagestyle{aftertitle}

%%%%%%%%%%%%%%%%%%%%%%%%%%%%%%%%%%%%%%%%%%%%%%%%%%%%%%%%%%%%
% === MAIN CONTENT ===
%%%%%%%%%%%%%%%%%%%%%%%%%%%%%%%%%%%%%%%%%%%%%%%%%%%%%%%%%%%%
\chapter{\DocTitle}
\section{مقدمه}

ارث‌بری، از پایه‌ای‌ترین مفاهیم شی‌گرایی است که به کلاس‌های مشابه، این امکان
را می‌دهد که متدها و فیلدهای یکسان را در یک \lr{class} جدید تعریف کنند.
با استفاده از ارث‌بری، می‌توانیم به طراحی‌ای بهتر برای برنامه‌مان برسیم به
طوری که کدهای تکراری کمی داشته باشیم و در نتیجه، برنامه
\lr{maintainable} تر باشد.

حتماً کدهای این داک رو اجرا کنید و هر سوالی که داشتید از تدریس‌یارها
بپرسید، در موردش سرچ کنید و از \lr{ChatGPT} راجع بهش اطلاعات بخواید.
ارث‌بری از مفاهیم مهم درس است که یادگیری آن اهمیت زیادی دارد و هسته‌ی اصلی
شی گرایی را تشکیل می‌دهد.

\section{ارث‌بری}

\subsection{توصیف یک مشکل}

در \lr{IntelliJ} یک پروژهٔ جدید ایجاد کنید و یک پکیج به اسم
\lr{\texttt{vehicle}} درست کنید. در \lr{\texttt{vehicle}} دو فایل جدید
به اسم \lr{\texttt{Car.java}} و \lr{\texttt{Bike.java}} درست کنید و داخل
اون‌ها، کلاس‌های \lr{\texttt{Car}} و \lr{\texttt{Bike}} رو ایجاد کنید:

\begin{codebox}{java}
package Vehicle;

public class Bike {
    public String brand;
    public int speed;
    public boolean hasCarrier;

    public Bike(String brand, int speed, boolean hasCarrier) {
        this.brand = brand;
        this.speed = speed;
        this.hasCarrier = hasCarrier;  
    }

    public void accelerate(int speed) {
        this.speed += speed;  
    }

    public void showBikeInfo() {  
        System.out.println("Brand: " + brand + ", Speed: " + speed + " km/h");  
        System.out.println("Has Carrier: " + hasCarrier);  
    }  
}
\end{codebox}

\begin{codebox}{java}
package Vehicle;

public class Car {
    public String brand;
    public int speed;
    public int doors;

    public Car(String brand, int speed, int doors) {
        this.brand = brand;
        this.speed = speed;
        this.doors = doors;  
    }

    public void accelerate(int speed) {
        this.speed += speed;  
    }

    public void showCarInfo() {  
        System.out.println("Brand: " + brand + ", Speed: " + speed + " km/h");  
        System.out.println("Doors: " + doors);  
    }  
}
\end{codebox}

بعد از این کار، کد \lr{\texttt{Main}} رو به کد زیر تغییر دهید و آن را
اجرا کنید تا مطمئن شوید که همه چیز درست کار می‌کند:

\begin{codebox}{java}
public static void main(String[] args) {
    Car myCar = new Car("Toyota", 180, 4);
    myCar.showCarInfo();

    System.out.println();

    Bike myBike = new Bike("Yamaha", 120, true);
    myBike.showBikeInfo();
}

\end{codebox}

بعد از اجرای کد، باید چنین چیزی ببینید:

\begin{codebox}{text}
Brand: Toyota, Speed: 180 km/h  
Doors: 4  
  
Brand: Yamaha, Speed: 120 km/h  
Has Carrier: true
\end{codebox}

حالا بیاید راجع به شباهت‌های این دو کلاس صحبت کنیم. ظاهر این کلاس‌ها خیلی
شبیه به هم‌ هستند. هر دو فیلدهایی به اسم \lr{\texttt{brand}} و
\lr{\texttt{speed}} دارند، با متدی به اسم \lr{\texttt{accelerate}} سرعت
وسیلهٔ نقلیه را زیاد می‌کنند و همینطور متدهایی دارند که اطلاعات آن‌ها را
نمایش می‌دهند. در واقع، اگر یک بار به کد نگاه کنید، می‌بینید کلاس‌های
\lr{\texttt{Bike}} و \lr{\texttt{Car}} تفاوت چندانی ندارند! حالا بیاید
یک کلاس جدید مثل \lr{\texttt{MotorCycle}} در فایل جدیدی به اسم
\lr{\texttt{MotorCycle.java}} بنویسیم:

\begin{codebox}{java}
package Vehicle;

public class MotorCycle {
    public String brand;
    public int speed;
    public boolean hasSidecar;

    public MotorCycle(String brand, int speed, boolean hasSidecar) {
        this.brand = brand;
        this.speed = speed;
        this.hasSidecar = hasSidecar;  
    }

    public void accelerate(int speed) {
        this.speed += speed;  
    }

    public void showMotorCycleInfo() {  
        System.out.println("Brand: " + brand + ", Speed: " + speed + " km/h");  
        System.out.println("Has Sidecar: " + hasSidecar);  
    }  
}
\end{codebox}

این کد را هم امتحان کنید:

\begin{codebox}{java}
MotorCycle myMotorCycle = new MotorCycle("Benz", 120, true);  
myMotorCycle.showBikeInfo();
\end{codebox}

باید خروجی‌ای مثل زیر ببینید:

\begin{codebox}{text}
Brand: Benz, Speed: 120 km/h  
Has Sidecar: true
\end{codebox}

کلاس جدیدی که ساختیم هم کار می‌کند، ولی باز هم شبیه‌ به دو کلاس قبلی است!
همانطور که می‌بینید هر وسیلهٔ نقلیه‌ای که بخواهیم اضافه کنیم، باید کدهای
کلاس‌های قبلی را تکرار کنیم. همانطور که در داک \lr{clean code} گفته شد،
ما در برنامه‌نویسی، سعی می‌کنیم از کدهای تکراری دوری کنیم.

\subsection{راه‌حل}

راه حل این مشکل این است که \textbf{ویژگی‌های مشترک تمام این کلاس‌ها را،
داخل کلاس جدیدی به اسم \lr{\texttt{Vehicle}} بنویسیم و بعد، با استفاده
از ارث‌بری، آن ویژگی‌ها را به همه این کلاس‌ها منتقل کنیم}. حالا این‌ که
«ویژگی‌های یک کلاس‌ را به کلاسی دیگر منتقل کنیم» یعنی چه؟ بیاید تا با کد
زدن ببینیم. توی فایل \lr{\texttt{Vehicle.java}} ، کلاس
\lr{\texttt{Vehicle}} رو ایجاد کنید. این کلاس، به جای این که نمایندهٔ
ماشین و موتور و دوچرخه باشد، نمایندهٔ یک «وسیلهٔ نقلیه‌» است و ویژگی‌های
مشترک تمام کلاس‌های قبل را داخل خود نگه می‌دارد:

\begin{codebox}{java}
package Vehicle;

public class Vehicle {
    public String brand;
    public int speed;

    public Vehicle(String brand, int speed) {
        this.brand = brand;
        this.speed = speed;  
    }

    public void accelerate(int speed) {
        this.speed += speed;  
    }

    public void showInfo() {  
        System.out.println("Brand: " + brand + ", Speed: " + speed + " km/h");  
    }  
}
\end{codebox}

می‌بینید که فیلدها و متدهای مشترک تمام وسایل نقلیه، یعنی
\lr{\texttt{brand}} ، \lr{\texttt{speed}} ، \lr{\texttt{accelerate}} و
غیره، همگی داخل این کلاس هستند. حالا بیاید کد کلاس \lr{\texttt{Car}} رو
عوض کنیم:

\begin{codebox}{java}
public class Car extends Vehicle {
    public int doors;

    public Car(String brand, int speed, int doors) {
        super(brand, speed);
        this.doors = doors;  
    }

    public void showCarInfo() {  
        showInfo();  
        System.out.println("Doors: " + doors);  
    }  
}
\end{codebox}

در این کد، شما برای اولین بار با کلیدواژهٔ \lr{\texttt{extends}} آشنا
می‌شوید. این کلیدواژه در جاوا برای تعیین ارث‌بری استفاده می‌شود. در کد بالا
کلاس \lr{\texttt{Car}} ، تمام فیلدها و متدهای کلاس \lr{\texttt{Vehicle}}
رو دارد. مثلا در خط اول متد \lr{\texttt{showCarInfo()}} ، می‌بینید که متد
\lr{\texttt{showInfo()}} صدا زده شده اما \lr{\texttt{Car}} چنین متدی
ندارد. در واقع، این متد در کلاس \lr{\texttt{Vehicle}} تعریف شده و چون
\lr{\texttt{Car}} از \lr{\texttt{Vehicle}} ارث‌بری کرده، پس می‌تواند از
متدها و فیلدهای آن استفاده بکند!

حالا در \lr{\texttt{main}} کد زیر را اجرا کنید تا مطمئن شوید همه چیز
درست کار می‌کند:

\begin{codebox}{java}
public static void main(String[] args) {  
    Car myCar = new Car("Toyota", 180, 4);  
    myCar.showCarInfo();  
  
    System.out.println();  
  
    myCar.accelerate(5);  
    myCar.showCarInfo();  
}
\end{codebox}

بعد از اجرای این کد، باید خروجی زیر را ببینید:

\begin{codebox}{text}
Brand: Toyota, Speed: 180 km/h  
Doors: 4  

Brand: Toyota, Speed: 185 km/h  
Doors: 4
\end{codebox}

می‌بینید که حتی توانستیم از متد \lr{\texttt{accelerate}} هم برای
\lr{\texttt{myCar}} استفاده کنیم، در حالی که این متد اصلا در
\lr{\texttt{Car}} وجود ندارد! این موضوع فقط به خاطر این است که
\lr{\texttt{Car}} از \lr{\texttt{Vehicle}} ارث‌بری می‌کند، و به همین دلیل
به تمام فیلدها و متدهای آن دسترسی دارد.

حالا کد‌ کلاس‌های \lr{\texttt{Bike}} و \lr{\texttt{MotorCycle}} را به شکل
زیر تغییر دهید تا آن‌ها هم از کلاس \lr{\texttt{Vehicle}} ارث‌بری کنند:

\begin{codebox}{java}
package Vehicle;

public class MotorCycle extends Vehicle {
    boolean hasSidecar;

    public MotorCycle(String brand, int speed, boolean hasSidecar) {
        super(brand, speed);
        this.hasSidecar = hasSidecar;  
    }

    public void showBikeInfo() {  
        showInfo();  
        System.out.println("Has Sidecar: " + hasSidecar);  
    }  
}
\end{codebox}

\begin{codebox}{java}
package Vehicle;

public class Bike extends Vehicle {
    public boolean hasCarrier;

    public Bike(String brand, int speed, boolean hasCarrier) {
        super(brand, speed);
        this.hasCarrier = hasCarrier;  
    }

    public void showMotorCycleInfo() {  
        showInfo();  
        System.out.println("Has Carrier: " + hasCarrier);  
    }  
}
\end{codebox}

می‌بینید که حجم کدهای تکراری خیلی کمتر شدند. با کد زیر، این کلاس‌های جدید
را تست کنید:

\begin{codebox}{java}
public static void main(String[] args) {  
    Bike myBike = new Bike("Yamaha", 120, true);  
    myBike.showBikeInfo();  
  
    System.out.println();  
      
    myBike.accelerate(5);  
    myBike.showBikeInfo();

    System.out.println();  
  
    MotorCycle myMotorcycle = new MotorCycle("Harley-Davidson", 160, false);  
    myMotorcycle.showMotorCycleInfo();  
}
\end{codebox}

باید خروجی زیر را ببینید:

\begin{codebox}{text}
Brand: Yamaha, Speed: 120 km/h  
Has Carrier: true  
  
Brand: Yamaha, Speed: 125 km/h  
Has Carrier: true  
  
Brand: Harley-Davidson, Speed: 160 km/h  
Has Sidecar: false
\end{codebox}

می‌بینید که همهٔ این کلاس‌ها، مثل کلاس \lr{Car} می‌توانند از فیلدهای
\lr{\texttt{brand}} و \lr{\texttt{speed}} و متدهای
\lr{\texttt{accelerate()}} و \lr{\texttt{showInfo()}} استفاده کنند.

\subsection{\lr{subclass} ها و \lr{superclass} ها}

در کد بالا، به کلاس \lr{\texttt{Vehicle}} که از آن ارث‌بری شده
\lr{superclass} یا کلاس والد می‌گویند. همچنین به کلاس‌های
\lr{\texttt{Car}} ، \lr{\texttt{Bike}} و \lr{\texttt{MotorCycle}} که از
کلاس \lr{\texttt{Vehicle}} ارث‌بری کردند \lr{subclass} یا کلاس‌های فرزند
می‌گویند. هر کلاس می‌تواند از حداکثر یک کلاس به صورت مستقیم ارث‌بری کند.
یعنی هر کلاس حداکثر یک والد دارد.

همینطور که می‌بینید در خط اول \lr{constructor} های \lr{subclass} ها، از
عبارتی به اسم \lr{\texttt{super}} استفاده شده است:

\begin{codebox}{java}
public Bike(String brand, int speed, boolean hasCarrier) {
    super(brand, speed);
    this.hasCarrier = hasCarrier;  
}
\end{codebox}

هر زمان از کلید واژه جدید \lr{super} استفاده می‌کنیم، یعنی قصد داریم از
فیلدها، متدها یا کانستراکتورهای کلاس پدر استفاده کنیم. مثلا در این‌جا، با
استفاده از کلیدواژهٔ \lr{\texttt{super}} ، نشان دادیم که می‌خواهیم
کانستراکتور زیر از کلاس پدر را استفاده کنیم تا فیلدهای
\lr{\texttt{brand}} و \lr{\texttt{speed}} برای آبجکت جدید مقداردهی شوند:

\begin{codebox}{java}
public Vehicle(String brand, int speed) {
    this.brand = brand;
    this.speed = speed;  
}
\end{codebox}

یا مثلا، در متد \lr{\texttt{showBikeInfo()}} ، می‌توانیم با استفاده از
\lr{\texttt{super}} ، دقیق‌تر مشخص کنیم که متد \lr{\texttt{showInfo}} ی
کلاس \lr{\texttt{Vehicle}} صدا زده شود:

\begin{codebox}{java}
public void showBikeInfo() {
    super.showInfo();  
    System.out.println("Has Carrier: " + hasCarrier);  
}
\end{codebox}

همچنین با استفاده از این کلیدواژه، شما می‌توانید به فیلدهای کلاس پدر هم
دسترسی داشته باشید:

\begin{codebox}{java}
public void showBikeInfo() {
    super.showInfo();  
    System.out.println("Has Carrier: " + hasCarrier);  
      
    System.out.println("Speed: " + super.speed);  
    System.out.println("Brand: " + super.brand);  
}
\end{codebox}

البته نتیجه کپی کردن کد بالا به جای کد \lr{\texttt{ShowBikeInfo()}} این
است که سرعت و برند دوچرخه‌ها دو بار چاپ می‌شوند! هدف از نشان دادن کد بالا
این بود که ببینید با کلیدواژهٔ \lr{super} ، به تمام ویژگی‌های کلاس پدر
دسترسی خواهید داشت. در ادامه این داک، متدهای کلاس پدر را \lr{Override}
می‌کنید و بیشتر به اهمیت این موضوع پی می‌برید.

\section{کلاس‌های \lr{Abstract}}

فرض کنید می‌خواهید یک برنامه بنویسید تا حقوق کارمندان را حساب کند. سازمان
شما دو نوع کارمند دارد:

\begin{itemize}
\tightlist
\item
  \textbf{کارمندهای تمام‌وقت:} این کارمندها، ماهیانه مبلغ ثابتی حقوق
  می‌گیرند. برای مثال ممکن است هر ماه ۳۰ میلیون تومان حقوق دریافت کنند.
\item
  \textbf{کارمندهای قراردادی:} این کارمندها، بر اساس ساعاتی که در ماه در
  شرکت بودند حقوق می‌گیرند. برای مثال اگر ساعتی ۱ میلیون تومان دریافت
  کنند ، اگر ۲۰ ساعت در شرکت باشند ۲۰ میلیون و اگر ۴۵ ساعت باشند، ۴۵
  میلیون حقوق دریافت می کنند.
\end{itemize}

پکیج \lr{\texttt{company}} را برای درست کردن کلاس‌ها ایجاد کنید. دوتا
کلاس \lr{\texttt{FullTimeEmployee}} و \lr{\texttt{Contractor}} را در این
پکیج ایجاد کنید:

\begin{codebox}{java}
package company;

public class Contractor {
    public String name;
    public int id;
    public double hourlyRate;
    public int hoursWorked;

    public Contractor(String name, int id, double hourlyRate, int hoursWorked) {
        this.name = name;
        this.id = id;
        this.hourlyRate = hourlyRate;
        this.hoursWorked = hoursWorked;  
    }

    public void showDetails() {  
        System.out.println("ID: " + id + ", Name: " + name);  
    }

    public double calculateSalary() {
        return hourlyRate * hoursWorked;  
    }  
}
\end{codebox}

\begin{codebox}{java}
package company;

public class FullTimeEmployee {
    public String name;
    public int id;
    public double monthlySalary;

    public FullTimeEmployee(String name, int id, double monthlySalary) {
        this.name = name;
        this.id = id;
        this.monthlySalary = monthlySalary;  
    }

    public void showDetails() {  
        System.out.println("ID: " + id + ", Name: " + name);  
    }

    public double calculateSalary() {
        return monthlySalary;  
    }  
}
\end{codebox}

در این دو کلاس، موارد مشترک زیادی می‌شود دید. مثل فیلدهای
\lr{\texttt{name}} و \lr{\texttt{id}} و متد \lr{\texttt{showDetails}} که
کاملا یکسان هستند. این شباهت، این حس را ایجاد می کند که باید از
\lr{inheritance} استفاده کنید. برای این کار، شما کلاس
\lr{\texttt{Employee}} رو تعریف می‌کنید:

\begin{codebox}{java}
package company;

public class Employee {
    public String name;
    public int id;

    public Employee(String name, int id) {
        this.name = name;
        this.id = id;  
    }

    public void showDetails() {  
        System.out.println("ID: " + id + ", Name: " + name);  
    }  
}
\end{codebox}

و کدتون رو به شکلی تغییر دهید که کلاس‌های \lr{\texttt{FullTimeEmployee}}
و \lr{\texttt{Contractor}} هر دو از آن ارث‌بری کند. کلاس‌های زیر را در
پکیج \lr{company} ایجاد کنید یا تغییر بدهید:

\begin{codebox}{java}
package company;

public class Contractor extends Employee {
    public double hourlyRate;
    public int hoursWorked;

    public Contractor(String name, int id, double hourlyRate, int hoursWorked) {
        super(name, id);
        this.hourlyRate = hourlyRate;
        this.hoursWorked = hoursWorked;  
    }

    public double calculateSalary() {
        return hourlyRate * hoursWorked;  
    }  
}
\end{codebox}

\begin{codebox}{java}
package company;

public class FullTimeEmployee extends Employee {
    public double monthlySalary;

    public FullTimeEmployee(String name, int id, double monthlySalary) {
        super(name, id);
        this.monthlySalary = monthlySalary;  
    }

    public double calculateSalary() {
        return monthlySalary;  
    }  
}
\end{codebox}

می‌بینید که کد تمیزتری شد. حالا بیاید در متد \lr{\texttt{main}} ، یک سری
کارمند تعریف کنیم و مجموع حقوقی که باید بدیم را حساب کنیم:

\begin{codebox}{java}
public static void main(String[] args) {  
    FullTimeEmployee ali = new FullTimeEmployee("Ali", 1, 10_000);  
    Contractor gholi = new Contractor("Gholi", 2, 1_000, 100);  
    Contractor mamad = new Contractor("Mamad", 3, 1_000, 25);  
  
    ArrayList<FullTimeEmployee> fullTimeEmployees = new ArrayList<>();  
    ArrayList<Contractor> contractors = new ArrayList<>();  
  
    fullTimeEmployees.add(ali);  
    contractors.add(mamad);  
    contractors.add(gholi);

    double sumOfSalary = 0;
    for (FullTimeEmployee employee : fullTimeEmployees) {  
        sumOfSalary += employee.calculateSalary();  
    }
    for (Contractor contractor : contractors) {  
        sumOfSalary += contractor.calculateSalary();  
    }  
      
    System.out.println("Salary this month: " + sumOfSalary);  
}
\end{codebox}

این کد را یکبار بخوانید تا بفهمید چه کاری انجام می‌دهد. آن را اجرا کنید
تا نتیجه زیر را ببینید:

\begin{codebox}{java}
Salary this month: 135000.0
\end{codebox}

یک مشکل در این کد وجود دارد. فرض کنید رئیس می‌خواهد که نوع جدیدی از
کارمند داشته باشیم. مثلا فرض کنید که می‌خواهیم حقوق مدیران به شکل متفاوتی
محاسبه شود. کلاس \lr{\texttt{Manager}} را اضافه می‌کنیم و بعد از این کار،
برای محاسبهٔ حقوق حلقهٔ زیر را به \lr{\texttt{main}} اضافه می‌کنیم:

\begin{codebox}{java}
for (Manager manager : managers) {  
    sumOfSalary += manager.calculateSalary();  
}
\end{codebox}

بعدتر تصمیم می گیرید که حقوق کله گنده‌ها به شکل متفاوتی محاسبه شود! شما
کلاس \lr{\texttt{KaleGonde}} را برای آن‌ها درست می‌کنید و برای حساب کردن
حقوق حلقهٔ زیر را اضافه می‌کنید:

\begin{codebox}{java}
for (KaleGonde kaleGonde : kaleGondeha) {  
    sumOfSalary += kaleGonde.calculateSalary();  
}
\end{codebox}

کد نهایی‌ چنین چیزی خواهد بود:

\begin{codebox}{java}
double sumOfSalary = 0;
for (FullTimeEmployee employee : fullTimeEmployees) {  
    sumOfSalary += employee.calculateSalary();  
}
for (Contractor contractor : contractors) {  
    sumOfSalary += contractor.calculateSalary();  
}
for (Manager manager : managers) {  
    sumOfSalary += manager.calculateSalary();  
}
for (KaleGonde kaleGonde : kaleGondeha) {  
    sumOfSalary += kaleGonde.calculateSalary();  
}
\end{codebox}

حدس بزنید مشکل این کد کجاست؟ بله دقیقا. این حلقه‌ها خیلی مشابه هستند و ما
کدهای شبیه به هم را ترجیح نمی‌دهیم! یک راه حل خوب برای این موضوع، این هست
که کلاس \lr{\texttt{Employee}} ، که پدر همهٔ کلاس‌های برنامه‌ است را
\lr{abstract} کنیم. در ابتدا کد را ببینیم:

\begin{codebox}{java}
package company;

public abstract class Employee {
    public String name;
    public int id;

    public Employee(String name, int id) {
        this.name = name;
        this.id = id;  
    }

    public void showDetails() {  
        System.out.println("ID: " + id + ", Name: " + name);  
    }

    public abstract double calculateSalary();  
}
\end{codebox}

با استفاده از کلیدواژهٔ \lr{abstract} در خط سوم این کد، کلاس
\lr{\texttt{Employee}} را به یک \lr{Abstract Class} تبدیل کردیم. می‌بینید
که در این کلاس، یک اتفاق عجیب افتاده. متد
\lr{\texttt{calculateSalary()}} تعریف شده است، ولی هیچ چیزی در
\lr{\texttt{\{\}}} برای پیاده‌سازی آن نیامده. این یعنی کلاس
\lr{\texttt{Employee}} ، پیاده‌سازی این متد را به \lr{subclass} های خود
واگذار کرده.

در واقع، کلاس \lr{\texttt{Employee}} به شما می گوید «هر کارمند، یک متد
به اسم \lr{\texttt{calculateSalary()}} دارد که هیچ چیزی ورودی نمی‌گیرد و
خروجی آن، \lr{\texttt{double}} خواهد بود. پیاده‌سازی این متد بسته به نوع
کارمند (\lr{\texttt{FullTime}} ، \lr{\texttt{Contractor}} ،
\lr{\texttt{Manager}} یا \lr{\texttt{KaleGonde}} بودن) متفاوت است، ولی
هر کلاس، قطعا این متد را دارد و شما می‌توانید آن را فراخوانی کنید.».

شاید مقداری گیج کننده باشد ، بیاید تا کد \lr{\texttt{main}} را یک مقدار
عوض کنیم تا این مطلب برایتان به خوبی مفهوم شود:

\begin{codebox}{java}
public static void main(String[] args) {  
    FullTimeEmployee ali = new FullTimeEmployee("Ali", 1, 10_000);  
    Contractor gholi = new Contractor("Gholi", 2, 1_000, 100);  
    Contractor mamad = new Contractor("Mamad", 3, 1_000, 25);  
  
    ArrayList<Employee> allEmployees = new ArrayList<>();  
  
    allEmployees.add(ali);  
    allEmployees.add(mamad);  
    allEmployees.add(gholi);

    double sumOfSalary = 0;
    for (Employee employee : allEmployees) {  
        sumOfSalary += employee.calculateSalary();  
    }  
  
    System.out.println("Salary this month: " + sumOfSalary);  
}
\end{codebox}

کد فعلی \lr{\texttt{main}} را نگاه کنید و تفاوت‌های آن را با کد قبلی
ببینید. می‌بینید که آرایه‌ ما از جنس \lr{\texttt{Employee}}است، و آن را با
\lr{\texttt{FullTimeEmployee}} ها و \lr{\texttt{Contractor}} ها پر
کردیم. از اون‌جایی که می‌دانیم هر \lr{\texttt{Employee}} فارغ از نوع قطعا
یک متد \lr{\texttt{calculateSalary()}} دارد، پس فقط یک حلقه کافیست که
درآمد همهٔ کارکنان را جمع بزنیم. اگر این کد اجرا شود، خروجی قبل را
می‌بینید:

\begin{codebox}{java}
Salary this month: 135000.0
\end{codebox}

کار \lr{abstract class} ها دقیقا همین چیزی است که این‌جا دیدید. وقتی در
کلاس پدر، نمی‌دانیم که کلاس‌های فرزند یک متد را به چه شکل پیاده‌سازی
می‌کنند، ولی می‌دانیم که حتماً آن را پیاده‌سازی می‌کنند، از
\lr{abstract class} استفاده می‌کنیم.

دقت کنید که دوتا نکته در \lr{abstract class} ها بسیار مهم است. اول این
که کلاس‌های فرزند، حتما باید متدهای \lr{abstract} ای که کلاس پدر تعریف
کرده پیاده‌سازی کنند، چرا که در غیر این صورت به خطا می‌خوریم. یک بار متد
\lr{\texttt{calculateSalary}} را از \lr{\texttt{Contractor}} پاک کنید تا
خطایش را ببینید:

\begin{figure}[H]
\centering
\includegraphics[width=0.8\textwidth]{images/oop-2-9.png}
\end{figure}

این خطا به شما می‌گوید که «کلاس \lr{\texttt{Contractor}} یا باید خودش
\lr{abstract} باشد (و در نتیجه پیاده‌سازی متدهای \lr{abstract} را به
فرزندانش بسپرد)، یا باید متد \lr{\texttt{calculateSalary}} را پیاده‌سازی
کند».

نکتهٔ دوم این است که چون کلاس‌های \lr{abstract} پیاده‌سازی کاملی ندارند،
شما نمی‌توانید از این کلاس‌ها مستقیما \lr{\texttt{new}} کنید. کد زیر به
خطا می‌خورد:

\begin{codebox}{java}
Employee javad = new Employee("Javad", 4);
\end{codebox}

کلاس‌های \lr{abstract} ، فقط ظرفی برای \lr{subclass} های خود هستند و
نباید از آنها مستقیم آبجکت ساخته بشود. در عوض، چیزی مثل کد زیر کاملا
درست است:

\begin{codebox}{java}
FullTimeEmployee ali = new FullTimeEmployee("Ali", 1, 10_000);  
Employee ali2 = ali;
\end{codebox}
	
	
	
	
%%%%%%%%%%%%%%%%%%%%%%%%%%%%%%%%%%%%%%%%%%%%%%%%%%%%%%%%%%%%
\end{document}